% General multilayer schematic of the deployed monotone core families.
% Ellipses denote repeated hidden layers; widths and curves are illustrative.
% Purple paths bundle nonnegative linear maps of the entire original input.
\begin{tikzpicture}[x=1cm,y=1cm,font=\small,
  weight/.style={draw=blue!70!black,line width=.65pt},
  connection/.style={draw=black!60,line width=.6pt},
  inputmap/.style={draw=violet!65!black,dashed,line width=.65pt},
  maparrow/.style={inputmap,-{Latex[length=1.4mm]}},
  dot/.style={circle,fill=black,inner sep=0pt,minimum size=4pt},
  callout/.style={-{Latex[length=1.8mm]},draw=black,line width=.75pt},
  annotation/.style={font=\small\itshape,align=center,text width=3.15cm,inner sep=2pt},
  pics/softplus/.style={code={
    \draw[draw=red!65!black,fill=white,line width=.45pt]
      (-.44,-.32) rectangle (.44,.32);
    \draw[black!25,line width=.2pt]
      (-.33,-.22) -- (.34,-.22) (-.33,-.22) -- (-.33,.24);
    \draw[red!70!black,line width=.7pt]
      plot[domain=-3:3,samples=35]
      ({.11*\x},{-.22+.42*(ln(1+exp(\x))-ln(1+exp(-3)))/3});
  }},
  pics/hatcurve/.style args={#1/#2}{code={
    \draw[draw=teal!70!black,fill=white,line width=.4pt]
      (-.44,-.32) rectangle (.44,.32);
    \draw[black!25,line width=.2pt]
      (-.33,-.22) -- (.34,-.22) (-.33,-.22) -- (-.33,.24);
    \draw[teal!75!black,line width=.6pt]
      plot[domain=-1.5:1.5,samples=41]
      ({.22*\x},{-.22+.42*(#1*(\x+1.5)
        +#2*(max(\x+1,0)^3-2*max(\x,0)^3+max(\x-1,0)^3)/6)
        /(3*#1+1.5*#2)});
  }}
]
 \path[use as bounding box] (0,0) rectangle (16.2,9.45);
 \fill[blue!3] (1.6,8.75) rectangle (8.9,9.45);
 \fill[teal!4] (8.9,8.75) rectangle (16.2,9.45);
 \fill[black!2] (0,1.1) rectangle (1.6,9.45);
 \draw[black!65,line width=.5pt] (0,0) rectangle (16.2,9.45);
 \foreach \tablex in {1.6,8.9}
   \draw[black!65,line width=.5pt] (\tablex,1.1) -- (\tablex,9.45);
 \foreach \tabley in {1.1,7.2,8.75}
   \draw[black!65,line width=.5pt] (0,\tabley) -- (16.2,\tabley);
 \node[font=\small\bfseries] at (.8,9.1) {Core};
 \node[font=\small\bfseries] at (5.25,9.1) {Monotone ICNN};
 \node[font=\small\bfseries] at (12.55,9.1) {Monotone integrated-hat ICKAN};
 \node[align=center,font=\footnotesize] at (.8,7.98) {Layer\\operation};
 \node[align=center,font=\footnotesize] at (.8,4.15) {Multilayer\\core};
 \node[align=center,text width=6.6cm] at (5.25,7.98)
   {\textcolor{blue!70!black}{Learned nonnegative linear weights} + biases\\[3pt]
    $\longrightarrow$ \textcolor{red!70!black}{fixed softplus at hidden nodes}};
 \node[align=center,text width=6.6cm] at (12.55,7.98)
   {\textcolor{teal!75!black}{One learned scalar function per edge}\\[2pt]
    \textcolor{teal!75!black}{(convex and nondecreasing)}\\[2pt]
    $\longrightarrow$ sum at nodes + biases};

 % ICNN: information flows from the conditioned input at the top to the
 % scalar output at the bottom, consistently with the paper overview.
 \begin{scope}[xshift=1.6cm,yshift=1.35cm]
  \node[anchor=west,font=\small\bfseries] at (.12,5.55) {(a)};
  \node[font=\footnotesize,text=violet!65!black] (icinputvector) at (1.6,5.62)
    {input $\bm y$};
  \foreach \level in {4.05,1.70}
    \filldraw[red!65!black,fill=red!3,line width=.75pt]
      (.02,\level-.25) rectangle (3.18,\level+.25);
  \node[dot] (icinput1) at (.9,5.10) {};
  \node[dot] (icinputm) at (2.3,5.10) {};
  \node[font=\footnotesize] at (.9,5.34) {$y_1$};
  \node[font=\footnotesize] at (2.3,5.34) {$y_m$};
  \node[font=\footnotesize] at (1.6,5.10) {$\cdots$};
  \node[dot] (icoutput) at (1.6,.55) {};
  \node[font=\footnotesize] at (1.6,.12) {$\widehat H(\bm y)$};
  \foreach \j in {0,...,2} {
    \pgfmathsetmacro{\hiddenx}{.35+1.25*\j}
    \draw[weight] (icinput1) -- (\hiddenx,4.05);
    \draw[weight] (icinputm) -- (\hiddenx,4.05);
    \draw[weight] (\hiddenx,1.70) -- (icoutput);
    \foreach \k in {0,...,2} {
      \pgfmathsetmacro{\nextx}{.35+1.25*\k}
      \draw[weight] (\hiddenx,4.05) -- (\nextx,1.70);
    }
  }
  \fill[white] (.02,2.70) rectangle (3.23,3.13);
  \foreach \hiddenx in {.35,1.6,2.85}
    \node at (\hiddenx,2.91) {$\vdots$};

  % Direct input map to the output, and bundled maps into later preactivations.
  \draw[inputmap,rounded corners=2pt] (icinputvector.east)
    -- (3.7,5.62) -- (3.7,.55);
  \draw[maparrow] (3.7,.55) -- (icoutput.east);
  \draw[inputmap] (3.7,2.22) -- (.35,2.22);
  \foreach \hiddenx in {.35,1.6,2.85}
    \draw[maparrow] (\hiddenx,2.22) -- (\hiddenx,1.95);
  \node[fill=white,inner sep=1pt,text=violet!65!black] at (3.7,2.91) {$\vdots$};
  \foreach \j in {0,...,2} {
    \pgfmathsetmacro{\hiddenx}{.35+1.25*\j}
    \pic[scale=.58] at (\hiddenx,4.05) {softplus};
    \pic[scale=.58] at (\hiddenx,1.70) {softplus};
  }
  \node[annotation,text=red!70!black] (icfixed) at (5.65,4.15)
    {Nonlinear, fixed\\softplus activations};
  \node[annotation,text=blue!70!black] (icweights) at (5.65,3.05)
    {Learned nonnegative\\linear weights};
  \node[annotation,text=violet!65!black] (icmaps) at (5.65,1.25)
    {Original input to later\\hidden layers and output};
  \draw[callout] (3.21,4.05) to[out=0,in=195] (icfixed.west);
  \draw[callout] (3.16,3.18) to[out=0,in=180] (icweights.west);
  \draw[callout] (3.77,1.25) -- (icmaps.west);
 \end{scope}

 % ICKAN: the same top-to-bottom convention; every displayed edge contains
 % its own learned scalar function before the receiving sum node.
 \begin{scope}[xshift=8.9cm,yshift=1.35cm]
  \node[anchor=west,font=\small\bfseries] at (.12,5.55) {(b)};
  \node[font=\footnotesize,text=violet!65!black] (kaninputvector) at (1.6,5.62)
    {input $\bm y$};
  \foreach \level in {4.55,2.25,1.05}
    \filldraw[teal!70!black,fill=teal!3,line width=.75pt]
      (0,\level-.22) rectangle (3.26,\level+.22);
  \node[dot] (kaninput1) at (.9,5.10) {};
  \node[dot] (kaninputm) at (2.3,5.10) {};
  \node[font=\footnotesize] at (.9,5.34) {$y_1$};
  \node[font=\footnotesize] at (2.3,5.34) {$y_m$};
  \node[font=\footnotesize] at (1.6,5.10) {$\cdots$};
  \node[dot] (kanoutput) at (1.6,.55) {};
  \node[font=\footnotesize] at (1.6,.12) {$\widehat H(\bm y)$};
  \foreach \j in {0,...,2} {
    \pgfmathsetmacro{\hiddenx}{.35+1.25*\j}
    \node[dot] (kanfirst\j) at (\hiddenx,3.95) {};
    \node[dot] (kanlast\j) at (\hiddenx,1.70) {};
  }
  \foreach \j in {0,...,2} {
    \pgfmathsetmacro{\hiddenx}{.35+1.25*\j}
    \draw[connection] (kanlast\j) -- (\hiddenx,1.05) -- (kanoutput);
    \foreach \r in {0,1} {
      \pgfmathsetmacro{\edgex}{.27+.53*(2*\j+\r)}
      \ifnum\r=0
        \draw[connection] (kaninput1) -- (\edgex,4.55) -- (kanfirst\j);
      \else
        \draw[connection] (kaninputm) -- (\edgex,4.55) -- (kanfirst\j);
      \fi
    }
    \foreach \k in {0,...,2} {
      \pgfmathsetmacro{\edgex}{.16+.36*(3*\j+\k)}
      \draw[connection] (kanfirst\k) -- (\edgex,2.25) -- (kanlast\j);
    }
  }
  \fill[white] (.02,2.78) rectangle (3.27,3.20);
  \foreach \hiddenx in {.35,1.6,2.85}
    \node at (\hiddenx,2.99) {$\vdots$};
  \foreach \j in {0,...,2} {
    \pgfmathsetmacro{\hiddenx}{.35+1.25*\j}
    \pgfmathsetmacro{\slope}{.08+.25*\j}
    \pgfmathsetmacro{\curvature}{1.5-.45*\j}
    \pic[scale=.45] at (\hiddenx,1.05) {hatcurve=\slope/\curvature};
    \foreach \r in {0,1} {
      \pgfmathsetmacro{\edgex}{.27+.53*(2*\j+\r)}
      \pgfmathsetmacro{\slope}{.08+.15*\j+.2*\r}
      \pgfmathsetmacro{\curvature}{1.4-.3*\j-.2*\r}
      \pic[scale=.39] at (\edgex,4.55) {hatcurve=\slope/\curvature};
    }
    \foreach \k in {0,...,2} {
      \pgfmathsetmacro{\edgex}{.16+.36*(3*\j+\k)}
      \pgfmathsetmacro{\slope}{.08+.11*\j+.14*\k}
      \pgfmathsetmacro{\curvature}{1.5-.23*\j-.25*\k}
      \pic[scale=.30] at (\edgex,2.25) {hatcurve=\slope/\curvature};
    }
  }
  \draw[inputmap,rounded corners=2pt] (kaninputvector.east)
    -- (3.7,5.62) -- (3.7,.55);
  \draw[maparrow] (3.7,.55) -- (kanoutput.east);
  \node[annotation,text=teal!75!black] (kanfunctions) at (5.65,4.42)
    {One learned convex,\\nondecreasing function\\per edge};
  \node[annotation] (kansums) at (5.65,3.13) {Sum operation\\at nodes};
  \node[annotation,text=violet!65!black] (kanmap) at (5.65,1.25)
    {Linear map from\\input to output};
  \draw[callout] (3.29,4.55) to[out=0,in=170] (kanfunctions.west);
  \draw[callout] (3.0,3.95) to[out=0,in=130] (kansums.west);
  \draw[callout] (3.77,1.25) -- (kanmap.west);
 \end{scope}

 % Colour legend and the sign constraints of these monotone implementations.
 \draw[weight,line width=1pt] (.35,.78) -- (.85,.78);
 \node[anchor=west,font=\footnotesize] at (.98,.78) {Learned linear weights};
 \pic[scale=.40] at (5.05,.78) {softplus};
 \node[anchor=west,font=\footnotesize] at (5.42,.78) {Fixed softplus};
 \pic[scale=.40] at (9.05,.78) {hatcurve=.1/1.3};
 \node[anchor=west,font=\footnotesize] at (9.42,.78) {Learned spline functions};
 \draw[inputmap,line width=1pt] (13.35,.78) -- (13.85,.78);
 \node[anchor=west,font=\footnotesize] at (13.98,.78) {Input maps};
 \node[font=\footnotesize] at (8.1,.29)
   {Glyph counts show operator placement, not model size; all linear weights are nonnegative.};
\end{tikzpicture}
